% front/translator-note.tex —— 译者说明
% 说明底本、翻译原则、术语与记号约定、排版与校勘处理。
% 注：每段必须写成一行（不得硬换行），否则 xeCJK 会把换行当成空格，句中留出空隙。
% 全文须控制在一页之内（约 40 行），超出一页会使第二页只剩寥寥数行。
\cleardoublepage
\phantomsection
\chapter*{译者说明}
\addcontentsline{toc}{chapter}{译者说明}
\markboth{译者说明}{译者说明}

\section*{一、关于底本}

本书据科学出版社 2015 年出版的 \textit{Numerical Linear Algebra and Its Applications} (Second Edition) 译出 (「信息与计算科学丛书」第 73 卷, ISBN 978-7-03-046425-5), 作者为金小庆, 魏益民, 赵志. 原书初版于 2004 年; 第二版订正了初版中的若干疏漏, 补入定理~\ref{thm:6-7} 的证明, 新增第 10 章 (空间分数阶扩散方程), 参考文献亦有扩充. 全书共 10 章, 依次讨论线性方程组的直接法与迭代法, Krylov 子空间方法, 非对称与对称特征值问题, 以及上述方法在常微分方程与分数阶扩散方程中的应用.

\section*{二、翻译原则}

\noindent\textbf{1.~忠实原文.} 译文以原书正文为准, 不增删命题, 公式, 例题与习题, 不改动定理的条件与结论; 章节划分与叙述次序均依原书. 原书未设脚注, 译文亦不加脚注.

\smallskip
\noindent\textbf{2.~术语.} 术语以国内通行译名为准, 如 Gaussian elimination 作「高斯消去法」; 首见时以「中文 (English)」附出英文原词, 并收入书末术语索引, 同一概念全书只用一种译名.

\smallskip
\noindent\textbf{3.~人名.} 数学家的人名 (Jacobi, Gauss, Householder, Krylov 等) 保留原文拼写, 不另行音译; 致谢中的人名照录其缩写形式.

\smallskip
\noindent\textbf{4.~文献.} 原书以 [1]--[100] 引用文献, 译文沿用原编号, 与书末「参考文献」一一对应.

\smallskip
\noindent\textbf{5.~编号.} 公式, 定理, 定义, 引理, 推论, 例一律按章编号 (如式 (3.1), 定理 3.1); 习题按章自 1 起连续编号, 均依原书体例.

\smallskip
\noindent\textbf{6.~记号.} 向量用粗正体 ($\mathbf{x}$), 矩阵用大写斜体 ($A$); 转置记作 $^{T}$, 共轭转置记作 $^{*}$; 不等式用 $\leqslant$, $\geqslant$; $\R$, $\C$ 为实数域与复数域; $\norm{\cdot}$ 为范数, 所指依上下文判定; 虚数单位记作 $\ii$.

\section*{三、排版与校勘}

译文以 \LaTeX{} 重排, 用 XeLaTeX 编译, 中文正文采用 Fandol 字体, 数学排版沿用 \texttt{amsmath}; 算法框, 示意图与全部公式均重新排印, 未使用扫描图像. 原书为影印本, 无文字层, 翻译时先以光学字符识别取得初稿, 再逐页比对页面图版核校, 凡定理陈述, 矩阵元素, 上下标与各类编号一律以图版为准; 个别疑似排印疏漏因无从确证作者原意, 译文一仍其旧.

\section*{四、关于译文}

译者力求译文准确, 术语统一, 行文可读, 惟学力有限, 疏漏恐所难免; 译文中的错误概由译者负责, 敬请读者指正. 原书的学术观点, 方法与结论均属原作者.

\vspace{0.8em}
\rightline{译者}
\rightline{2026 年 11 月}
